\documentclass[11pt]{article}
\usepackage{../../latex/stat4ml_style}

\title{Chapter 2: Moments \\ \large Part I --- Statistics Foundation for ML}
\author{Omid Safarzadeh}
\date{}

\begin{document}
\maketitle

\begin{abstract}
\noindent
A distribution is an infinite-dimensional object; moments are the finite,
computable summaries we actually work with -- the mean, the variance, the
covariance between two variables, and, packaged together, the moment
generating function. Every loss function you will minimize later in this
course (mean squared error, the Gaussian log-likelihood, the bias--variance
decomposition) is a statement about moments. This chapter builds them
carefully, in scalar and then matrix form, because the matrix form is
exactly the shape you will need for multivariate regression and for the
covariance matrices behind PCA (Part II, Chapter 5).

\emph{Acknowledgement: this chapter expands the STAT4ML lecture slides
(``Slide 2 -- Moments''), following Casella \& Berger, \emph{Statistical
Inference}, 2nd ed. (2002).}
\end{abstract}

\tableofcontents

\section{Moments}

\begin{definition}
For integer $n$, the \emph{$n$th moment} of $X$ is $\mu_n' = E[X^n]$. The
\emph{$n$th central moment} of $X$ is $\mu_n = E[(X-\mu)^n]$, where
$\mu = \mu_1' = E[X]$.
\end{definition}

\subsection{Expected Value}

An everyday ``average'' weights every observed data point equally. The
\emph{expected value} instead weights every \emph{possible} value a random
variable can take by how probable that value is. It answers: on average,
which value do we expect $X$ to take?

\begin{definition}
The expected value of $g(X)$ is
\[
  E[g(X)] =
  \begin{cases}
    \displaystyle\int_{-\infty}^{\infty} g(x)\, f_X(x)\, dx & X \text{ continuous} \\[1em]
    \displaystyle\sum_{x \in \mathcal{X}} g(x)\, f_X(x) & X \text{ discrete,}
  \end{cases}
\]
provided the sum/integral converges absolutely; if $E[|g(X)|] = \infty$ we
say $E[g(X)]$ does not exist.
\end{definition}

\begin{example}[Exponential mean]
Let $X \sim \mathrm{Exponential}(\lambda)$, $f_X(x) = \frac{1}{\lambda}
e^{-x/\lambda}$ for $x \ge 0$. Integrating by parts ($\int u\,dv = uv - \int
v\,du$, with $u = x$, $dv = \frac{1}{\lambda}e^{-x/\lambda}dx$),
\[
  E[X] = \int_0^\infty x \cdot \frac{1}{\lambda} e^{-x/\lambda}\,dx
       = \big[-xe^{-x/\lambda}\big]_0^\infty + \int_0^\infty e^{-x/\lambda}\,dx
       = \lambda.
\]
\end{example}

Expectation is a \emph{linear} operator: for constants $a, b$,
\[
  E[a + bX] = a + b\,E[X] = a + b\mu.
\]

\begin{theorem}
Let $a, b, c$ be constants and $g_1, g_2$ functions whose expectations exist.
\begin{enumerate}
  \item $E[a\,g_1(X) + b\,g_2(X) + c] = a\,E[g_1(X)] + b\,E[g_2(X)] + c$.
  \item If $g_1(x) \ge 0$ for all $x$, then $E[g_1(X)] \ge 0$.
  \item If $g_1(x) \ge g_2(x)$ for all $x$, then $E[g_1(X)] \ge E[g_2(X)]$.
  \item If $a \le g_1(x) \le b$ for all $x$, then $a \le E[g_1(X)] \le b$.
\end{enumerate}
\end{theorem}

\begin{example}
Let $X \sim \mathrm{Uniform}(0,1)$ and $g(X) = -\log X$. Then, using
integration by parts,
\[
  E[-\log X] = \int_0^1 -\log x\,dx = \big[-x\log x + x\big]_0^1 = 1.
\]
\end{example}

\subsection{Variance}

Variance measures the spread of $X$ around its mean. While the mean is
usually written $\mu$, variance is written $\sigma^2$; it is a second-order
moment.

\begin{definition}
The variance of $X$ is its second central moment,
\[
  \mathrm{Var}(X) = E[(X-\mu)^2], \qquad
  \sqrt{\mathrm{Var}(X)} = \text{standard deviation of } X.
\]
Equivalently, expanding the square and using linearity of expectation,
\[
  \mathrm{Var}(X) = E[X^2] - \mu^2.
\]
\end{definition}

This second form, $E[X^2] - (E[X])^2$, is the one you will actually compute
in code -- it needs only two running sums (of $x$ and of $x^2$) instead of a
second pass over the data. It is also the single most common source of bugs
in hand-written variance code: forgetting the $-\mu^2$ term always returns a
number too large, and never negative, which makes the bug easy to miss on
a quick sanity check. (This exact bug is planted in
\texttt{exercise02\_buggy.py} -- watch for it.)

\section{Covariance and Correlation}

\begin{definition}
For random variables $X, Y$ with $E[X] = \mu_X$, $E[Y] = \mu_Y$,
$\mathrm{Var}(X) = \sigma_X^2$, $\mathrm{Var}(Y) = \sigma_Y^2$, the
\emph{covariance} is
\[
  \mathrm{Cov}(X,Y) = E\big[(X-\mu_X)(Y-\mu_Y)\big],
\]
and the \emph{correlation} (correlation coefficient) is
\[
  \rho_{XY} = \frac{\mathrm{Cov}(X,Y)}{\sigma_X \sigma_Y}.
\]
\end{definition}

If large values of $X$ tend to co-occur with large values of $Y$ (and small
with small), $\mathrm{Cov}(X,Y) > 0$: when $X > \mu_X$, $Y > \mu_Y$ is also
likely, so $(X-\mu_X)(Y-\mu_Y) > 0$ on average. Symmetrically, if large
values of $X$ tend to pair with \emph{small} values of $Y$,
$\mathrm{Cov}(X,Y) < 0$. Correlation just rescales covariance by both
standard deviations so that it is unitless and bounded in $[-1, 1]$ --
which is what makes it comparable across different pairs of variables.

\begin{theorem}
$\mathrm{Cov}(X,Y) = E[XY] - \mu_X \mu_Y$.
\end{theorem}

\begin{theorem}
If $X \perp\!\!\!\perp Y$ (independent), then $\mathrm{Cov}(X,Y) =
\rho_{XY} = 0$. The converse is \textbf{false} in general: zero correlation
does not imply independence (a standard counterexample: $X \sim
\mathrm{Uniform}(-1,1)$, $Y = X^2$; then $\mathrm{Cov}(X,Y) = 0$ but $Y$ is a
deterministic function of $X$).
\end{theorem}

\begin{theorem}
For random variables $X, Y$ and constants $a, b$,
\[
  \mathrm{Var}(aX + bY) = a^2 \mathrm{Var}(X) + b^2 \mathrm{Var}(Y) +
  2ab\,\mathrm{Cov}(X,Y),
\]
which reduces to $a^2\mathrm{Var}(X) + b^2\mathrm{Var}(Y)$ when $X, Y$ are
independent (since $\mathrm{Cov}(X,Y)=0$).
\end{theorem}

If $X, Y$ are positively correlated, $\mathrm{Var}(X+Y) >
\mathrm{Var}(X)+\mathrm{Var}(Y)$ (the variances reinforce each other); if
negatively correlated, $\mathrm{Var}(X+Y) < \mathrm{Var}(X)+\mathrm{Var}(Y)$
(they partially cancel).

\subsection{Variance of Sums of Random Variables}
For constants $a_i$ and random variables $X_i$, $i=1,\dots,n$,
\[
  \mathrm{Var}\Big(\sum_{i=1}^n a_i X_i\Big)
  = \sum_{i=1}^n a_i^2 \mathrm{Var}(X_i) + \sum_{i \ne j} a_i a_j
  \mathrm{Cov}(X_i, X_j).
\]

\subsection{Higher Moments}
The third and fourth (central, standardized) moments -- skewness and
kurtosis -- describe how asymmetric and how fat-tailed a distribution is,
respectively. They matter in ML whenever a model's error distribution is
assumed Gaussian but the real residuals are heavy-tailed (e.g.\ financial
returns, or gradient noise in SGD with a small batch size).

\section{Moment Generating Functions}

\begin{definition}
The \emph{moment generating function} (MGF) of $X$ is
\[
  M_X(t) = E[e^{tX}],
\]
provided this expectation is finite for $t$ in some open interval containing
0 (i.e.\ there is $h>0$ such that $E[e^{tX}]$ exists for all $-h<t<h$). If no
such neighborhood exists, the MGF does not exist.
\end{definition}

\begin{theorem}
If $X$ has MGF $M_X(t)$, then $E[X^n] = M_X^{(n)}(0)$, the $n$th derivative
of $M_X$ evaluated at $t=0$. This is exactly what ``moment generating''
means: differentiate the MGF $n$ times at zero and you get the $n$th raw
moment, all at once.
\end{theorem}

\subsection{The Normal MGF}
For $X \sim N(\mu, \sigma^2)$ with $f(x) = \frac{1}{\sqrt{2\pi\sigma^2}}
\exp\big[-\tfrac{1}{2}\big(\tfrac{x-\mu}{\sigma}\big)^2\big]$, completing the
square in the exponent of $E[e^{tX}]$ gives
\[
  M_X(t) = \exp\Big(\mu t + \frac{\sigma^2 t^2}{2}\Big).
\]
Differentiating,
\[
  E[X] = M_X'(0) = \mu, \qquad
  E[X^2] = M_X''(0) = \sigma^2 + \mu^2, \qquad
  \mathrm{Var}(X) = E[X^2] - (E[X])^2 = \sigma^2.
\]
This is a good closed-form to numerically double-check: Exercise 2 asks you
to estimate the Normal MGF by direct simulation (average $e^{tX}$ over many
samples) and compare against the formula above.

\section{Matrix Notation for Moments}

Let $X$ be $(r \times 1)$ and $Y$ be $(c \times 1)$ random vectors. The
\emph{cross-covariance matrix} is
\[
  \mathrm{Cov}(X,Y) =
  \begin{pmatrix}
    \mathrm{Cov}(X_1,Y_1) & \cdots & \mathrm{Cov}(X_1,Y_c) \\
    \vdots & \ddots & \vdots \\
    \mathrm{Cov}(X_r,Y_1) & \cdots & \mathrm{Cov}(X_r,Y_c)
  \end{pmatrix}
  = E\big[(X - E[X])(Y-E[Y])^\top\big].
\]

For a single $(c\times 1)$ vector $X$, write $\mathrm{Cov}(X)$ for
$\mathrm{Cov}(X,X)$:
\[
  \mathrm{Cov}(X) =
  \begin{pmatrix}
    \mathrm{Var}(X_1) & \cdots & \mathrm{Cov}(X_1,X_c) \\
    \vdots & \ddots & \vdots \\
    \mathrm{Cov}(X_1,X_c) & \cdots & \mathrm{Var}(X_c)
  \end{pmatrix},
\]
a symmetric $(c \times c)$ matrix -- \emph{the} object PCA diagonalizes in
Part II, Chapter 5.

For a block vector $X = \binom{Y}{Z}$ with $Y$ $(p\times 1)$ and $Z$
$(q\times 1)$,
\[
  \mathrm{Cov}(X) =
  \begin{pmatrix}
    \mathrm{Cov}(Y) & \mathrm{Cov}(Y,Z) \\
    \mathrm{Cov}(Z,Y) & \mathrm{Cov}(Z)
  \end{pmatrix}.
\]

For non-stochastic vectors $a$ $(r\times1)$, $b$ $(c\times1)$, and
$\Sigma_{XY} = \mathrm{Cov}(X,Y)$,
\[
  \mathrm{Cov}(a^\top X, b^\top Y) = a^\top \Sigma_{XY} b.
\]
In particular, with $\Sigma = \mathrm{Var}(X)$,
\[
  \mathrm{Var}(a^\top X) = a^\top \Sigma a,
\]
and for two random vectors of the same dimension,
\[
  \mathrm{Var}(X+Y) = \Sigma_{XX} + \Sigma_{XY} + \Sigma_{YX} + \Sigma_{YY},
  \qquad
  \mathrm{Var}(a^\top(X+Y)) = a^\top\Sigma_{XX}a + 2a^\top\Sigma_{XY}a +
  a^\top\Sigma_{YY}a
\]
(using $a^\top \Sigma_{XY} a = a^\top \Sigma_{YX} a$, since both equal the
scalar $\mathrm{Cov}(a^\top X, a^\top Y)$).

Finally, these results extend from vectors to matrices: for a
non-stochastic matrix $R$,
\[
  E[RX] = R\,E[X], \qquad
  \mathrm{Var}(RX) = R\,\mathrm{Var}(X)\,R^\top = R\,\Sigma_{XX}\,R^\top.
\]
This is the identity that tells you how a linear layer $z = Wx$ transforms
the covariance of its input -- the starting point for reasoning about
activation statistics in a deep network.

\section{Summary}
\begin{summarybox}
\begin{itemize}[nosep]
  \item Moments ($E[X^n]$, central moments $E[(X-\mu)^n]$) are finite
        numerical summaries of a distribution.
  \item Variance and covariance are the second-moment objects that drive
        the bias--variance tradeoff (Part II, Chapter 2) and PCA (Part II,
        Chapter 5).
  \item The MGF packages every moment into one function via
        $E[X^n] = M_X^{(n)}(0)$.
  \item In matrix form, $\mathrm{Var}(RX) = R\Sigma R^\top$ is the identity
        you will use repeatedly once random vectors, not just random
        scalars, are on the table.
\end{itemize}
\end{summarybox}

\section{Exercises}
See \texttt{exercises/python/exercise02.py}, \texttt{exercises/r/exercise02.R},
and \texttt{exercises/rust/exercise02.rs} (expected value, variance,
covariance, correlation, and a Monte-Carlo check of the Normal MGF), plus the
Broken counterparts to debug in \texttt{exercise02\_buggy.\{py,R,rs\}}.

\section*{Reference}
Casella, G., \& Berger, R. (2002). \emph{Statistical Inference} (2nd ed.).
Cengage Learning.

\end{document}
